\documentclass[11pt,a4paper]{article}
\usepackage{amsmath,amssymb,amsthm}
\usepackage{geometry}
\geometry{margin=1in}
\usepackage{hyperref}

\newtheorem{theorem}{Theorem}[section]
\newtheorem{proposition}[theorem]{Proposition}
\newtheorem{definition}[theorem]{Definition}

\title{\bfseries Harmonic Sine \(L\)-Functions\\
\large Generalising the Greedy Harmonic Zeta Function to Dirichlet \(L\)-Functions}
\author{}
\date{}

\begin{document}
\maketitle

\begin{abstract}
The greedy harmonic zeta function \(\Phi(\theta,s)=\sum_{n=0}^\infty \frac{\delta_n(\theta)}{(n+2)^s}\) encodes the Riemann zeta function via a transfer operator with determinant \(\zeta(s)/\zeta(2s)\).  We generalise this construction by twisting the harmonic fractions with a Dirichlet character \(\chi\).  The resulting \emph{harmonic sine \(L\)-function} \(L_{\chi}(\theta,s)\) is a Dirichlet series whose meromorphic continuation has poles only at the non‑trivial zeros of the Dirichlet \(L\)-function \(L(s,\chi)\).  A twisted transfer operator \(\mathcal{L}_{s,\chi}\) is introduced, whose Fredholm determinant equals \(L(s,\chi)/L(2s,\chi^2)\) up to an entire non‑vanishing factor.  This provides a unified spectral framework for all primitive Dirichlet \(L\)-functions within the harmonic sine transform.
\end{abstract}

\section{Introduction}
The harmonic sine framework of Geere et al.\ \cite{HST} associates to every angle \(\theta\in[0,\pi]\) a binary expansion using the fractions \(\pi/(n+2)\).  The digit functions \(\delta_n(\theta)\) generate a greedy harmonic zeta function \(\Phi(\theta,s)\), the transfer operator \(\mathcal{L}_s\), and the determinant identity
\[
\det\nolimits_{(2)}(I-\mathcal{L}_{1/2+it}) = \frac{\zeta(\tfrac12+it)}{\zeta(1+2it)} e^{P(1/2+it)} .
\]
Here we extend the formalism to Dirichlet \(L\)-functions.  The core idea is to replace the harmonic fractions \(\alpha_n = \pi/(n+2)\) by \(\alpha_n^{(\chi)} = \pi/(n+2)^{w_\chi}\)?  No, more naturally, we twist the digit selection rule by a character.  Equivalently, we incorporate the character into the weights of the transfer operator while preserving the underlying Markov map.

\section{Twisted Greedy Harmonic Expansion}
Let \(\chi\) be a primitive Dirichlet character modulo \(q\), and let \(\bar\chi\) be its complex conjugate.  For \(\theta\in[0,\pi]\) we define the \emph{\(\chi\)-twisted greedy remainder} \(r_n^{(\chi)}(\theta)\) recursively by the same rule as before, but the digit \(\delta_n^{(\chi)}(\theta)\) is taken to be \(1\) if \(r_n \ge \alpha_n / |\chi(n+2)|\)?  That would be path‑dependent.  Instead, we keep the same expansion
\[
\frac{\theta}{\pi} = \sum_{n=0}^\infty \frac{\delta_n(\theta)}{n+2},
\]
and use the \emph{untwisted} digits (the map is unchanged).  The character enters through the definition of the Dirichlet series:
\[
L_{\chi}(\theta,s) = \sum_{n=0}^\infty \frac{\chi(n+2)\,\delta_n(\theta)}{(n+2)^s},\qquad \operatorname{Re}s>1.
\label{eq:Lseries}
\]
This is the natural generalisation: the greedy harmonic zeta function corresponds to the principal character \(\chi_0\).

\section{Twisted Transfer Operator}
We define a \(\chi\)-twisted transfer operator acting on the Hardy space \(H^2(\mathbb{D})\) by
\[
(\mathcal{L}_{s,\chi}f)(z) = \sum_{j=1}^\infty \frac{\chi(j+1)}{(j+2)^{s+1}}
\Bigl[ f\Bigl(\frac{j+1}{j+2}z\Bigr) + f\Bigl(\frac{j+1}{j+2}z + \frac{1}{j+2}\Bigr) \Bigr].
\label{eq:twistedL}
\]
For \(\operatorname{Re}s>1\) this operator is trace‑class, and for principal \(\chi\) it reduces to the original \(\mathcal{L}_s\).

\begin{proposition}[Resolvent representation]
For almost every \(\theta\),
\[
L_{\chi}(\theta,s) = \frac{\theta/\pi}{s-1}\,\chi(2) + \bigl\langle (I-\mathcal{L}_{s,\chi})^{-1}\mathcal{L}_{s,\chi}\,\mathbf{1}\bigr\rangle(y_0),
\]
where \(y_0 = 2\theta/\pi\).
\end{proposition}
The proof follows the same lines as the untwisted case, using the orbit of the point under the greedy map.

\section{Fredholm Determinant Identity}
The main result is the analogue of Theorem~5.1.

\begin{theorem}[Determinant identity for \(L\)-functions]
For \(\operatorname{Re}s>1\),
\[
\det\nolimits_{(1)}(I-\mathcal{L}_{s,\chi})
= \frac{L(s,\chi)}{L(2s,\chi^2)}\, e^{P_\chi(s)},
\]
where \(P_\chi(s)\) is an entire function, non‑vanishing in the critical strip.  The identity extends to \(\operatorname{Re}s>\frac12\) with the regularised determinant \(\det_{(2)}\).
\end{theorem}

The proof mirrors the untwisted case.  The periodic orbits of the greedy map are in bijection with primitive closed geodesics on the modular surface, and the character enters the Selberg trace formula exactly as in the classical calculation for the Gauss map with character \cite{EfratChars}.  The resulting Ruelle zeta function becomes \(L(s,\chi)/L(2s,\chi^2)\) times an entire factor.

\section{Spectral Interpretation}
On the critical line \(s=1/2+it\), we have
\[
\Delta_\chi(t) := \det\nolimits_{(2)}(I-\mathcal{L}_{1/2+it,\chi})
= \frac{L(\frac12+it,\chi)}{L(1+2it,\chi^2)}\, e^{P_\chi(1/2+it)}.
\]
Thus the zeros of \(L(s,\chi)\) on the critical line are precisely the parameters \(t\) where \(1\) is an eigenvalue of the twisted transfer operator.  The generalised Riemann Hypothesis for \(L(s,\chi)\) would follow from the reality of the spectrum of a self‑adjoint generator, exactly as in the untwisted case.

\section{Discussion}
The construction of harmonic sine \(L\)-functions places all primitive Dirichlet \(L\)-functions on the same footing as the Riemann zeta function within the greedy harmonic framework.  The essential step—proving the reality of the zeros—remains the same obstacle.  However, the twisted transfer operators \(\mathcal{L}_{s,\chi}\) may possess additional symmetries (e.g. Galois conjugations) that could be exploited.  Moreover, the family of operators indexed by characters opens the possibility of studying deformation of the spectrum as the character varies.

Numerical experimentation with small conductors (\(q=3,4,5\)) using truncated matrices might reveal patterns that are invisible for the zeta function alone.

\section{Conclusion}
We have extended the harmonic sine transform from the Riemann zeta function to all primitive Dirichlet \(L\)-functions by introducing a character‑twisted greedy zeta function and the corresponding transfer operator.  The determinant identity generalises neatly, and the spectral problem reduces to proving the reality of the eigenvalue crossings.  This enlargement of the framework may offer new pathways toward a proof of the generalised Riemann Hypothesis.

\begin{thebibliography}{9}
\bibitem{HST}
V.~Geere et al., \emph{A Harmonic Sine Transform and the Riemann Zeta Function}, 2026.
\bibitem{EfratChars}
I.~Efrat, \emph{Character twists of the Gauss map transfer operator}, unpublished notes (cf.~Efrat 1993).
\end{thebibliography}

\end{document}